\chapter{\label{chap:discussion}Discussion}

% This chapter must be completed from the frozen offline evidence and the
% supervised live records. It must not convert mocked success into provider
% evidence or generalize beyond the reviewed Small scenarios.

\section{Answer to RQ1: Operationalization}

The implemented method operationalizes the theoretical model through a chain
of versioned and digest-linked boundaries: immutable workload, provider
capability profiles, pricing evidence, cost ledger, Resolved Twin Architecture,
Resolved Deployment Specification, Deployment Manifest, graph-derived
readiness, durable operation, telemetry verification, and cleanup evidence.
This demonstrates the engineering path from a theoretical cost assignment to
an executable and inspectable deployment decision.

The final answer must distinguish three levels of evidence:

\begin{enumerate}
    \item deterministic contract and integration evidence;
    \item provider readiness, identity, and image-probe evidence; and
    \item supervised Apply, telemetry, Destroy, inventory, and observed-cost
    evidence.
\end{enumerate}

Any provider blocker found during the live phase remains part of the answer. It
identifies an operationalization boundary and must not be replaced by an
offline fixture.

\section{Answer to RQ2: Functional Comparability}

The comparison unit is a curated provider service bundle rather than a product
name. AWS, Azure, and GCP may satisfy one logical responsibility through
different combinations of managed services, hosted software, adapters,
identity mechanisms, and operations. A candidate enters the cost comparison
only when all mandatory components and edges resolve to explicit deployment,
verification, cleanup, and pricing ownership.

The final discussion should evaluate where the shared contract provides a
meaningful comparison and where provider-specific semantics remain. In
particular, it should discuss delivery and replay behavior, the hosted GCP Twin
and Grafana bundles, workload-identity differences, regional availability,
quota behavior, and the limits of treating a curated implementation as
representative of an entire provider.

\section{Answer to RQ3: Cost Effect}

The quantitative answer must use the frozen workload and price evidence and
report estimated monthly cost rather than observed billing unless provider
billing data is separately reconciled. For each baseline and multi-cloud
candidate, the discussion should report:

\begin{itemize}
    \item the exact component assignment and admissibility result;
    \item total estimated monthly cost and difference from each provider-local
    baseline;
    \item component, transfer, bridge, and shared-pool contributions;
    \item price, workload, tier, and region sensitivity; and
    \item the difference between the estimate and any observed short-lived
    evaluation cost.
\end{itemize}

No conclusion about a cost advantage should be written until the final
calculation digests and the supervised evidence records have been frozen.

\subsection{RQ3.1: Baseline Control}

The three provider-local baselines and every selected multi-cloud candidate
must use the same Small workload, Six-layer profile, currency, pricing policy,
regions, and functional requirements. This paired design isolates provider
allocation from changes in scenario semantics. The nine live cases cover all
three local providers and every directed AWS/Azure/GCP pair without claiming
to enumerate every valid component permutation.

\subsection{RQ3.2: Explicit Eventing Responsibility}

The explicit Eventing responsibility should be assessed in three dimensions:
functional behavior, topology, and cost. The discussion must trace ordering,
delivery, retry, replay, terminal-failure handling, and directed bridge
identity, then attribute the corresponding Eventing and bridge costs. This
does not claim that Eventing universally improves a Digital Twin; it evaluates
one bounded, versioned contract and workload family.

\section{Architecture and Pattern Trade-offs}

The retained Strategy and provider-profile patterns demonstrate that cost
scoring and provider implementations can be extended without exposing inactive
objectives or runtime plugins. Only the cost strategy and the three reviewed
provider profiles are executable. This is a deliberate compromise: the code
keeps clear seams, while the PoC avoids promising a generic optimization or
architecture framework.

The immutable Twin lifecycle similarly trades operational flexibility for
traceability. Duplicate, Import, Deploy, Destroy, and post-Destroy Re-deploy
are sufficient for the evaluation. In-place Terraform updates, migrations,
rollbacks, and continuous re-optimization would require new evidence and state
contracts and are therefore outside scope.

Durable SSE replay adds implementation complexity but protects a meaningful
safety property: reconnecting the UI cannot duplicate a cost-incurring Apply
or Destroy. In contrast, application OAuth, SAML, RBAC, and multi-tenancy would
not strengthen the research answers and were removed from the runtime path.

\section{Simplifications and Limitations}

The principal limitations are:

\begin{itemize}
    \item one local owner profile and no production authentication or
    multi-tenancy;
    \item one executable Six-layer profile, three regions, and one curated
    provider bundle per responsibility;
    \item monetary cost as the only optimization objective;
    \item frozen pricing snapshots rather than continuous price management or
    billing reconciliation;
    \item immutable deployed Twins and no in-place infrastructure update;
    \item bounded Python user functions rather than arbitrary executable
    deployment projects;
    \item provider account creation, billing recovery, quota approval,
    organization policy, and broad consent remain external;
    \item no embedded dashboard administration or general operations console;
    and
    \item nine supervised Small scenarios rather than exhaustive architecture,
    region, workload, concurrency, or failure-injection coverage.
\end{itemize}

\section{Threats to Validity}

\subsection{Internal Validity}

Implementation defects, transient provider failures, stale readiness results,
and incomplete cleanup can affect a live run. Immutable digests, one active
scenario, explicit confirmations, operation replay, immediate Destroy, and
post-Destroy inventory reduce these risks. They do not eliminate them. Every
deviation must remain attached to its scenario record.

\subsection{Construct Validity}

Functional completeness is defined by the chosen contract. It cannot prove
that differently composed provider bundles are equivalent for every Digital
Twin use case. Likewise, estimated cost depends on normalized workload and
pricing assumptions. Sensitivity analysis and explicit exclusions are needed
to avoid turning a modeling decision into a universal claim.

\subsection{External Validity}

The results are limited to AWS, Azure, and GCP, the selected regions and
service versions, the Small evaluation workload, and the fixed Six-layer
contract. They do not automatically generalize to other providers, regions,
Digital Twin topologies, production load, concurrent users, or long-running
operations.

\subsection{Conclusion Validity}

Small cost differences may be dominated by pricing drift, minimum billing
units, short-lived provisioning effects, or unmodeled operational effort. The
final conclusion must therefore report absolute and relative differences,
sensitivity, observed deviations, and uncertainty rather than relying only on
the identity of the cheapest candidate.
